% ============================================================
% Section 5: Discussion
% ============================================================

\section{Discussion}
\label{sec:discussion}

% \section{考察}
% \label{sec:discussion}
%
% 本章では，評価結果を踏まえて，プリエンプティブ方式が唯一機能する理由，
% チェックポイントオーバーヘッドのトレードオフ，および実用上の示唆を考察する．

% ============================================================
\subsection{Why Preemptive Is the Only Sustainable Policy}
\label{sec:discussion_why}
% ============================================================

% \subsection{なぜプリエンプティブ方式だけが持続可能なのか}
% \label{sec:discussion_why}

The simulation results reveal a structural difference between the three
sharing policies that goes beyond a matter of degree.
Under FCFS and Owner Priority, there exists a range of high system load
($\rho_\mathrm{sys} \geq 0.8$--$0.9$) in which Tier9 owners observe PR $> 1$
and therefore have a strictly rational incentive to withdraw from ACP.
Once they withdraw, the high-performance GPUs disappear from the shared pool,
degrading the utility of sharing for all remaining users.
This triggers a participation cascade (Section~\ref{sec:sustainability})
that ultimately collapses the entire system.

% 評価結果から，3つの共有方式には単なる程度の差ではなく，構造的な違いがあることが示された．
% FCFSおよび所有者優先方式では，高負荷域（$\rho_\mathrm{sys} \geq 0.8$〜$0.9$）において
% Tier9所有者がPR $> 1$を観測し，ACPからの離脱が合理的となる負荷帯が存在する．
% 離脱によって高性能GPUが共有プールから消えると，残る全ユーザーの共有効用が低下し，
% 参加カスケード崩壊（\ref{sec:sustainability}節）を引き起こして
% 最終的にシステム全体が崩壊する．

Preemptive avoids this cascade by construction.
When an owner's task arrives, the guest job is suspended immediately,
and the owner's GPU is returned with zero queuing delay.
The owner's TAT under ACP is therefore at most equal to the standalone TAT,
i.e., PR $\leq 1$ holds by design, not just empirically.
This provides a \emph{structural} participation guarantee:
regardless of how many other users are sharing the GPU,
the owner never faces a worse outcome than not participating.
As a result, no rational Tier9 owner ever has an incentive to withdraw,
the high-performance GPUs remain in the pool permanently,
and the system converges to a stable, full-participation equilibrium.

% プリエンプティブ方式は設計上このカスケードを回避する．
% 所有者タスクが到着した瞬間にゲストジョブが即時停止され，
% GPUはキュー待機時間ゼロで所有者に返却される．
% ACP下での所有者TATはしたがって，単独利用時のTATと高々等しい，
% すなわちPR $\leq 1$が経験的にではなく設計上保証される．
% これは\emph{構造的}な参加保証を提供する：
% 何人の他ユーザーがGPUを共有していても，所有者は不参加より悪い結果に直面しない．
% その結果，合理的なTier9所有者が離脱するインセンティブを持つことはなく，
% 高性能GPUは常時共有プールに留まり，システムは安定した全参加平衡へ収束する．

The key insight is that \emph{preemption serves as ownership insurance,
not merely a throughput optimization}.
Existing literature on GPU preemption (e.g., Chimera~\cite{han2023chimera},
AntMan~\cite{xiao2020antman}) treats preemption as a mechanism to improve
aggregate resource utilization.
In ACP, preemption is reinterpreted as a mechanism to guarantee that
sharing never harms the owner---converting ACP participation from a
gamble into a dominated strategy.

% 重要な洞察は，\emph{プリエンプションがスループット最適化に留まらず，
% 所有権保険として機能する}という点である．
% GPU プリエンプションに関する既存研究（Chimera~\cite{han2023chimera}，
% AntMan~\cite{xiao2020antman}等）では，プリエンプションは
% 資源利用率を向上させる効率化ツールとして位置づけられている．
% ACPでは，プリエンプションは「共有が所有者を絶対に損させない」ことを保証する
% メカニズムとして再解釈される――ACP参加を賭博から支配戦略へと転換する．

% ============================================================
\subsection{The Checkpoint Overhead Trade-off}
\label{sec:discussion_overhead}
% ============================================================

% \subsection{チェックポイントオーバーヘッドのトレードオフ}
% \label{sec:discussion_overhead}

The Preemptive policy imposes a checkpoint-and-resume penalty of factor
$\alpha = 0.2$ on every evicted guest task.
This overhead has two effects that must be considered separately.

% プリエンプティブ方式は，退去する全ゲストタスクにチェックポイント・再開ペナルティ
% 係数$\alpha = 0.2$を課す．このオーバーヘッドは2種類の影響をもたらし，それぞれを分けて考える必要がある．

\textbf{Effect on guests.}
Evicted guest tasks accumulate overhead across multiple preemptions,
particularly when running on a highly contested GPU (A100).
At $\rho_\mathrm{sys} \geq 0.9$, frequent evictions raise the effective
task size of guest jobs, contributing to overall system TAT growth.
This is a real cost, but it is borne by non-owner users---the same users
who benefit from access to high-performance resources the rest of the time.
Since ACP is not a throughput-maximization system but a participation-based
sharing ecosystem, this trade-off is consistent with the design goal.

% \textbf{ゲストへの影響．}
% 退去ゲストタスクは，特に競合の激しいGPU（A100）で実行されている場合，
% 複数回のプリエンプションにわたってオーバーヘッドが蓄積する．
% $\rho_\mathrm{sys} \geq 0.9$では頻繁な退去がゲストジョブの実効タスクサイズを増大させ，
% システム全体のTAT上昇に寄与する．
% これは現実のコストだが，非所有者ユーザー（他の時間帯に高性能資源へのアクセス恩恵を受ける
% 同一ユーザー）が負担する．
% ACPはスループット最大化システムではなく参加型共有エコシステムであるため，
% このトレードオフは設計目標と整合する．

\textbf{Effect on owners.}
For the owner, the checkpoint overhead is irrelevant in most scenarios:
the owner triggers preemption only when submitting their own task,
and their task starts immediately without waiting for the checkpoint to complete
(the checkpoint runs in parallel with, or prior to, the owner's task start,
depending on implementation).
The exception identified in the evaluation is the Low-Heavy scenario at
$\rho_\mathrm{sys} \geq 0.9$: when the owner's inference task is
\emph{shorter} than $\alpha \times$ (size of the interrupted training job),
the preemption latency can marginally inflate the owner's TAT.
This is an inherent limitation of checkpoint-based preemption for
inference-dominant owners, and motivates future work on
\emph{kill-and-restart} or \emph{migration-free preemption} strategies
(Section~\ref{sec:future_work}).

% \textbf{所有者への影響．}
% 所有者にとって，チェックポイントオーバーヘッドはほとんどのシナリオで無関係である．
% 所有者は自身のタスク投入時にのみプリエンプションを発動し，
% タスクはチェックポイント完了を待たずに開始される
% （実装に応じてチェックポイントは並行または事前に実行される）．
% 評価で特定された例外は，$\rho_\mathrm{sys} \geq 0.9$のLow-Heavyシナリオである：
% 所有者の推論タスクが$\alpha \times$（中断された学習ジョブのサイズ）より
% \emph{短い}場合，プリエンプションレイテンシが所有者TATをわずかに膨張させることがある．
% これはチェックポイント型プリエンプションの本質的な限界であり，
% 将来の\emph{kill-and-restart}や\emph{migration-free preemption}戦略への動機となる
% （\ref{sec:future_work}節）．

% ============================================================
\subsection{Relation to the Incentive Model}
\label{sec:discussion_incentive}
% ============================================================

% \subsection{インセンティブモデルとの関係}
% \label{sec:discussion_incentive}

The decision model (Section~\ref{sec:model}) defines participation as
rational when $U_i \geq 1$, which is equivalent to PR $\leq 1$.
The simulation results confirm that Preemptive maintains this condition
at all load levels and in three out of four heterogeneous scenarios,
providing empirical support for the theoretical guarantee.

% 意思決定モデル（\ref{sec:model}節）は，$U_i \geq 1$のとき参加が合理的であると定義し，
% これはPR $\leq 1$と等価である．
% シミュレーション結果は，プリエンプティブ方式が全負荷・4シナリオ中3つで
% この条件を維持することを確認し，理論的保証の経験的裏付けを提供する．

The iterative participation model further shows that the equilibrium
is not merely \emph{stable} but \emph{attractive}: even starting from
a random initial state where fewer than half the users participate,
the system converges to full participation within 2--3 iterations
under Preemptive.
This \emph{convergence property} is critical for practical deployment:
ACP does not require all users to join simultaneously or to trust each
other in advance---rational individual decisions alone lead to the
socially optimal outcome.

% 反復型参加モデルはさらに，平衡が単に\emph{安定}なだけでなく\emph{吸引的}であることを示す：
% 半数未満のユーザーが参加するランダムな初期状態から始まっても，
% プリエンプティブ方式では2〜3イテレーションで全参加状態へ収束する．
% この\emph{収束特性}は実用展開において重要である：
% ACPは全ユーザーの同時参加や事前の相互信頼を必要とせず，
% 合理的な個人の意思決定だけで社会的最適な結果がもたらされる．

% ============================================================
\subsection{Practical Deployment Considerations}
\label{sec:discussion_practical}
% ============================================================

% \subsection{実用展開上の考察}
% \label{sec:discussion_practical}

\textbf{Checkpoint and migration cost.}
The $\alpha = 0.2$ penalty used in this study is a conservative estimate
consistent with GPU checkpoint overhead reported for training workloads
in the literature~\cite{xiao2020antman,han2023chimera}.
For inference tasks, which typically have much smaller memory footprints,
$\alpha$ may be substantially lower.
Calibrating $\alpha$ for specific GPU models and workload types is
a necessary step before production deployment.

% \textbf{チェックポイント・マイグレーションコスト．}
% 本研究で用いた$\alpha = 0.2$は，学習ワークロードにおける
% GPUチェックポイントオーバーヘッドの文献値と整合する保守的な推定値である
% \cite{xiao2020antman,han2023chimera}．
% メモリフットプリントが通常はるかに小さい推論タスクでは，$\alpha$はかなり低くなる可能性がある．
% 本番環境へのデプロイ前に，特定のGPUモデルとワークロード種別に対して
% $\alpha$を較正することが必要なステップである．

\textbf{Scale.}
The 18-user, 9-tier configuration is representative of a small-to-medium
campus research environment.
Scaling to hundreds of users would require a distributed scheduler
capable of tracking GPU ownership and preempting jobs across nodes.
The core scheduling logic (preempt on owner arrival) is simple and
stateless per GPU, making it amenable to distributed implementation.
We expect the PR $\leq 1$ guarantee to hold at scale,
since it depends only on the per-GPU owner-preemption rule,
not on global coordination.

% \textbf{スケール．}
% 18ユーザー・9ティアの構成は，小〜中規模の学内研究環境を代表する．
% 数百ユーザーへのスケールアップには，GPU所有権を追跡しノードをまたいで
% ジョブをプリエンプトできる分散スケジューラが必要となる．
% コアとなるスケジューリングロジック（所有者到着時のプリエンプト）は
% GPU単位でシンプルかつステートレスであるため，分散実装に適している．
% PR $\leq 1$保証はGPU単位の所有者プリエンプションルールにのみ依存し，
% グローバル協調を必要としないため，スケール時も成立すると期待される．

\textbf{Integration with existing schedulers.}
Preemptive Owner-Priority scheduling can be implemented as an extension
to existing GPU cluster schedulers such as Kubernetes with GPU plugins
or SLURM with GRES support.
The primary addition required is an ownership table (mapping GPU $\to$ owner)
and a preemption trigger invoked when an owner's job is submitted.
The resume strategy (restart on own GPU, wait, or migrate) can be
configured per workload type to minimize the effective $\alpha$.

% \textbf{既存スケジューラへの統合．}
% プリエンプティブ所有者優先スケジューリングは，GPUプラグイン付きKubernetesや
% GRESサポート付きSLURMなどの既存GPUクラスタスケジューラへの拡張として実装できる．
% 主に必要な追加は，所有権テーブル（GPU→所有者のマッピング）と，
% 所有者ジョブ投入時に呼び出されるプリエンプショントリガーである．
% 再開戦略（自GPU再起動・待機・マイグレーション）は，実効$\alpha$を最小化するために
% ワークロード種別ごとに設定可能である．
